% *****************************************************************************
% **************            SLIDE
% *****************************************************************************
\beamerdefaultoverlayspecification{}
\begin{frame}[fragile]                           % IL FAUT METTRE FRAGILE A CHAQUE FOIS QU'IL Y A DU CODE A METTRE :(
\frametitle{Un exemple de code}

\vspace*{-1mm}
\begin{block}{Un petit code}                    % Hélas, \visible<1->{   NE FONCTIONNE PAS :(   (pb de compatibilité entre package ...)
\vspace*{-2mm}
\begin{footnotesize}
\begin{java}
public class {
|	...
|	public static void main(String arg){
|	|	System.out.println("sdmklfjsdfj dfsdf ");
|	|	for(int i=0;i<n;i++)
|	|	|	if (1=1)
|	|	|	|	Toto.add(i)
|	}
}
\end{java}
\end{footnotesize}
\vspace*{-2mm}
\end{block}


\vspace*{-1mm}

\visible<2->{
\begin{block}{Commentaires provenant du dossier ShareText}
\begin{itemize}
	\item Citer vos sources~;
	\item Référencer les figures dans le texte~;
	\item Avoir une notation cohérente entre les différents supports~;
	\item Respecter le typographie pour le bien être du relecteur.
\end{itemize}

\end{block}
}

\end{frame}
\beamerdefaultoverlayspecification{<+->}



% *****************************************************************************
% **************            SLIDE
% *****************************************************************************
\beamerdefaultoverlayspecification{}
\begin{frame}[fragile]                           % IL FAUT METTRE FRAGILE A CHAQUE FOIS QU'IL Y A DU CODE A METTRE :(
\frametitle{Un peu de maths}

\visible<1->{
\begin{block}{La formule dans \textcolor{redUPEC}{ShareText}}
\input{../ShareText/maths1.tex}
\end{block}
}

\visible<2->{
\begin{block}{Le tableau dans \textcolor{redUPEC}{ShareText}}
\begin{small}
\begin{tabular}{|l|l|l|}
\hline
Département & Nombre de chaises & Nombre de prises électriques\\
\hline
\hline
ITS & 150 & 20\\
\hline
SI & 149 & 25 \\
\hline
ISBS & 151 & 19\\
\hline
\end{tabular}
\end{small}
\end{block}
}

\visible<3->{
\begin{alertblock}{Théoreme, un résultat très important !}
\[
1+(-1)^n=\begin{cases}
			0, & \text{if $n$ odd}\\
            2, & \text{otherwise}
		 \end{cases}
\]
\end{alertblock}
}

\end{frame}
\beamerdefaultoverlayspecification{<+->}